\chapter{Social Change, Development, Globalization \& Work and Economy}

\section{Conceptual Issues \& Sociological Theories of Social Change}

\subsection{Oxford Mobility Study 1972 (Goldthorpe)}

\begin{DataFact}
\begin{itemize}
    \item \textbf{Goldthorpe's second study} (after \textit{Affluent Workers}, 1963) is the \textbf{Oxford Mobility Study (1972)} --- the aim was to test whether \textbf{Britain} is an open or closed society. He surveyed \textbf{10,000 males in England and Wales} and divided them into \textbf{three clusters}: (1) the \textbf{service class} (highly privileged), (2) the \textbf{intermediate class}, and (3) the \textbf{working class}.
    \item \textbf{Finding 1 --- high absolute mobility}: with the growth of the \textbf{service sector} (IT/software, hospitality --- better pay and life chances), everyone achieved upward, absolute, intra-generational mobility.
    \item \textbf{Finding 2 --- relative mobility unchanged}: those \textbf{born higher up the social ladder} always had \textbf{better chances} of reaching higher positions than those born lower. The \textbf{odds weigh more in favour of those from higher classes} --- so \textbf{equality of opportunity was not achieved}, and he concluded that \textbf{Britain is highly closed}.
    \item \textbf{Exam tip}: just mentioning \textbf{``Goldthorpe''} and \textbf{``Oxford Mobility Study''} (year optional) conveys the point --- society may \textit{look} open, but only the already-privileged few reap the benefits.
\end{itemize}
\end{DataFact}

\begin{QuickRevision}
\begin{itemize}
    \item Oxford Mobility Study (Goldthorpe, 1972): 10,000 males, 3 clusters (service/intermediate/working); high \textit{absolute} but unchanged \textit{relative} mobility; Britain highly closed.
\end{itemize}
\end{QuickRevision}

\sectiondivider

\subsection{Conceptual Issues: Equality, Hierarchy, Inequality, Differences, Stratification}

\begin{KeyConcept}
\begin{itemize}
    \item \textbf{Equality of opportunity} (all those aged 21--32 may appear for UPSC, tested on one objective paper) vs \textbf{equality of outcome} (everyone clears). \textbf{Marx} focuses on \textbf{equality of outcome} --- communism/classless society, but he left communism under-elaborated.
    \item Our constitution promises \textbf{equality of opportunity, never equality of outcome}. Merton: opportunity to reach the goal is equally available, but the \textbf{means are not} --- hence \textbf{deviance}.
    \item \textbf{Weber}: even under communism there would be no equality --- power and prestige remain (even a tribal society has a chief who alone wears the leopard skin). So \textbf{equality is a utopia; inequality is an ideology} --- we have all accepted inequality (if exclusion is ``based on merit,'' nobody revolts).
\end{itemize}
\end{KeyConcept}

\begin{KeyConcept}
\begin{itemize}
    \item \textbf{Hierarchy} = vertical arrangement; a ranked pattern (there can be no ``horizontal hierarchy''). It is one feature of bureaucracy (Weber's chain of command).
    \item \textbf{Hierarchy and inequality are synonyms} --- hierarchy implies inequality, inequality rests on hierarchical arrangement. But they are \textbf{not always based on merit} --- \textbf{Murphy} says racism/sexism/classism are justified in the name of meritocracy. (The IQ test justified white supremacy and European hegemony; Einstein got more fame than Bose for the Einstein--Bose state of matter because he was white/European --- Merton's comparative advantage.)
    \item Murphy's example: UPSC gives \textbf{equal opportunity to Hindi and English medium alike}, but in the name of merit the \textbf{English-medium candidate is promoted as an IAS officer} --- even though Hindi-medium people are the majority. Many prejudices reproduce themselves \textit{in the name of merit}.
\end{itemize}
\end{KeyConcept}

\begin{KeyConcept}
\begin{itemize}
    \item \textbf{Differences that cannot be quantified} (caste, religion, gender, language --- unlike class, which is quantifiable) are still \textbf{put into hierarchy}. Brahmin vs Kshatriya are only \textit{different}, like apple and orange --- you cannot say one is superior.
    \item \textbf{Dipankar Gupta}: differences \textbf{cannot assume hierarchical relation}, but in the \textbf{popular consciousness they do} (mainly due to \textbf{power structures} --- English is superior because the West colonised; Brahmins are superior because they wrote the Arthashastra and designed the caste system; men because rulers are male). \textbf{Class (quantifiable) $\rightarrow$ open strata; caste/gender/religion (subjective) $\rightarrow$ closed strata.}
    \item Even ``Brahmin'' is not homogeneous --- within it are Mishra, Pandey, Varma, Sharma, Trivedi, Chaturvedi, Dwivedi, Lingayat, all claiming superiority over one another (horizontally placed, claiming vertical relation). Likewise all \textbf{Kshatriya princely families} (Madhya Pradesh, Rajasthan, Gujarat) belong to ``one world'' yet each claims its own lineage/descent superior to the others.
    \item Why do Hindus assume superiority over Muslims in India? Partly \textbf{institutional communalism} --- Hindus being the numerical majority have more \textbf{social capital}, hence more \textbf{symbolic capital}; but that is not the only reason.
\end{itemize}
\end{KeyConcept}

\begin{KeyConcept}
\begin{itemize}
    \item \textbf{Sorokin} introduced the concept: \textbf{stratification = vertical arrangement of groups according to the distribution of power, wealth, income or prestige.}
    \item \textbf{Dipankar Gupta critiqued and modified} Sorokin: caste is \textbf{not necessarily vertical} (there are horizontal layers within a caste). He added \textbf{social differences that cannot be objectively measured but assume inequality}.
    \item If we study only differences $\rightarrow$ \textbf{social differentiation}; only hierarchy $\rightarrow$ \textbf{social inequality}; both together $\rightarrow$ \textbf{social stratification} (the larger umbrella). \textbf{Strata are not only vertical, they are also horizontal.}
\end{itemize}
\end{KeyConcept}

\begin{QuickRevision}
\begin{itemize}
    \item Equality of opportunity vs outcome (Marx = outcome); equality = utopia, inequality = ideology.
    \item Hierarchy = vertical ranked pattern; hierarchy/inequality = synonyms, not always merit-based (Murphy).
    \item Dipankar Gupta: unquantifiable differences assume hierarchy via power structures; class = open, caste/gender = closed; Brahmin not homogeneous.
    \item Sorokin's stratification; Gupta's modification; differentiation (differences) vs inequality (hierarchy) vs stratification (both).
\end{itemize}
\end{QuickRevision}

\sectiondivider

\subsection{Open \& Closed Systems}

\begin{ComparisonBox}
\begin{itemize}
    \item \textbf{Open system} = hierarchy is \textbf{quantifiable and legitimised on meritocratic grounds}; mobility \textbf{allowed via merit}; upward mobility \textbf{does not displace} the upwardly situated member (an SDM's promotion does not remove the existing IAS); mobility is of the \textbf{individual}; movement is with \textbf{relative ease}; it draws sustenance from \textbf{quantitative} hierarchy. Example: \textbf{class}.
    \item \textbf{Closed system} = differences based on race/caste/estate/gender are \textbf{hierarchised without objective criteria}; mobility is \textbf{discouraged} by the upper ones (whites resist black mobility; men resist women's education); upward mobility \textbf{displaces} a superior position of another group (if a Shudra becomes Brahmin, Brahmins lose status); mobility is of the \textbf{group} (movements); movement is \textbf{resisted}; it draws from \textbf{qualitative} differences. Example: \textbf{caste, estate, race, gender}.
    \item These are \textbf{ideal types}, not absolute reality --- neither is class absolutely open, nor caste absolutely closed.
\end{itemize}
\end{ComparisonBox}

\begin{KeyConcept}
\begin{itemize}
    \item In the closed system, upward mobility always displaces someone --- status is a \textbf{zero-sum game} (Murray Milner): one's gain is another's loss (hence the 50\% OBC ceiling; those already OBC resist newcomers).
\end{itemize}
\end{KeyConcept}

\begin{KeyConcept}
\begin{itemize}
    \item Movements (feminism, black/civil rights, Dalit/Dalit Panther, Republican Party of India, self-respect movement) arise in \textbf{closed} societies because mobility is resisted; open societies (class, merit) need no movement.
\end{itemize}
\end{KeyConcept}

\begin{CriticalPoint}
\begin{itemize}
    \item \textbf{Siddharth Gautam} (a Kshatriya) became \textbf{Lord Buddha} --- mobility to Brahminical status. \textbf{Valmiki} (a dalit) wrote the Ramayana and became \textbf{Maharshi Valmiki}. \textbf{Guru Nanak} (a Hindu) attained enlightenment and became the first Sikh Guru (Brahminical status). These are \textbf{individual} exceptions (sanyasis, gurus, sadhus) noted even by \textbf{Louis Dumont} (Homo Hierarchicus).
    \item Marriage: an upper-caste girl marrying a lower-caste boy $\rightarrow$ both \textbf{outcast} (the child a ``chandal'') = downward mobility --- proof it is open.
    \item \textbf{Westernisation}: cumulative $\rightarrow$ dispersed inequality (caste no longer defines occupation). \textbf{Reservation}: people who were ``general'' became OBC (1993); caste is now a \textbf{springboard} for mobility. \textbf{Sanskritisation} (M. N. Srinivas: lower castes emulate upper castes; the Coorg tribe became ``Coorg Brahmin''), and its reverse \textbf{de-Sanskritisation/Sudra-isation} (Brahmins becoming ``Shudra'' for reservation; Jats/Gujjars, dominant Kshatriya castes, claiming OBC; the Marathas of Maharashtra).
\end{itemize}
\end{CriticalPoint}

\begin{QuickRevision}
\begin{itemize}
    \item Open = quantifiable + meritocratic hierarchy, individual mobility, no displacement; Closed = subjective hierarchy, resisted group mobility, zero-sum displacement.
    \item Murray Milner (zero-sum status); movements = sign of closed society; caste not absolutely closed (Buddha, Valmiki, Nanak, reservation, Sanskritisation).
\end{itemize}
\end{QuickRevision}

\sectiondivider

\subsection{Evolutionary Theory of Social Change}

\begin{KeyConcept}
\begin{itemize}
    \item \textbf{Auguste Comte}: theological $\rightarrow$ metaphysical $\rightarrow$ scientific.
    \item \textbf{James Frazer}: magic $\rightarrow$ religion $\rightarrow$ science.
    \item \textbf{E. B. Tylor}: animism $\rightarrow$ polytheism $\rightarrow$ monotheism.
    \item \textbf{Herbert Spencer} (a Darwin follower, a \textbf{social Darwinist} --- Darwin spoke of species, Spencer extended evolution to societies): \textbf{simple $\rightarrow$ doubly compound $\rightarrow$ trebly compound} society. Multiple families form a \textbf{clan} (imagined descent --- nagavanshi, chandravanshi, \textbf{gotra} = descent from the \textbf{sapta rishi}); clans form a \textbf{tribe}; tribes form a \textbf{nation}.
    \item Traditional Hindu marriage checks \textbf{caste endogamy, gotra exogamy, village exogamy}.
\end{itemize}
\end{KeyConcept}

\begin{CriticalPoint}
\begin{itemize}
    \item Evolutionary theory implies that \textbf{progress = development}, and that societies progress in \textbf{stages} (you cannot skip a stage) --- like biology (you become an adult only through stages).
\end{itemize}
\end{CriticalPoint}

\begin{QuickRevision}
\begin{itemize}
    \item Comte (theological$\rightarrow$metaphysical$\rightarrow$scientific), Frazer (magic$\rightarrow$religion$\rightarrow$science), Tylor (animism$\rightarrow$polytheism$\rightarrow$monotheism), Spencer (social Darwinist; simple$\rightarrow$doubly$\rightarrow$trebly compound).
    \item Implies progress = development, stage-wise, unskippable.
\end{itemize}
\end{QuickRevision}

\sectiondivider

\subsection{Functionalist Theory (Talcott Parsons)}

\begin{KeyConcept}
\begin{itemize}
    \item \textbf{Talcott Parsons} gave the functionalist theory of social change. Imagine a man walking a \textbf{tight rope} with weights in both hands: he was in \textbf{equilibrium}, it is \textbf{disturbed} while walking, and \textbf{restored} when he finishes. Societies work the same way.
    \item Sources of change: \textbf{demography, technology, ideas}. When change comes, society \textbf{adapts} (structural adjustment, adaptive capacity) and restores equilibrium --- e.g., over/under-population $\rightarrow$ state policy; technological change $\rightarrow$ society adapts.
    \item Population control: \textbf{educate women} (a ``natural contraceptive'') or \textbf{instil fear} (China's one-child policy; India's Emergency). Coercive top-down measures show weak adaptive capacity; \textbf{Michel Foucault} calls this \textbf{biopower} (power over bodies). When the state regulates population from the top, \textbf{civil society challenges it from below} so that the state is forced towards \textbf{soft measures} --- this is how equilibrium is restored.
\end{itemize}
\end{KeyConcept}

\begin{ExampleBox}
\begin{itemize}
    \item \textbf{Demonetisation} (cashless transactions, Jan Dhan-Aadhaar-Mobile ``trinity'', Aadhaar-enabled payment) caused chaos but society maintained itself; families adapted to lockdown through apps (husbands helped women).
    \item \textbf{Lord Buddha} challenging Brahmanism and \textbf{Martin Luther} challenging Catholicism disturbed equilibrium; then society restored it (Buddhism/Hinduism now coexist; Protestantism arose).
\end{itemize}
\end{ExampleBox}

\begin{CriticalPoint}
\begin{itemize}
    \item Parsons talks of \textbf{moving equilibrium}, never revolution --- social change is the consequence of continuous structural adjustment (structural differentiation).
\end{itemize}
\end{CriticalPoint}

\begin{QuickRevision}
\begin{itemize}
    \item Parsons: equilibrium $\rightarrow$ disturbance $\rightarrow$ restoration (tight-rope); sources = demography/technology/ideas; structural adjustment + adaptive capacity.
    \item Population control (educate women, fear); Foucault's biopower; no revolution, only moving equilibrium.
\end{itemize}
\end{QuickRevision}

\sectiondivider

\subsection{Cultural Lag (William Ogburn)}

\begin{KeyConcept}
\begin{itemize}
    \item \textbf{William Ogburn} refined Parsons. Change happens at two levels: \textbf{material (technology --- rapid)} and \textbf{non-material (values --- slow)}. iPhone 12 $\rightarrow$ 13 changes fast; values change slowly.
    \item When material change outruns non-material change, there is \textbf{cultural lag} --- our technology is modern but our values are not. This causes \textbf{strain} in the system.
    \item \textbf{Luddites} are the people who resist technological innovation in the name of ethics/values.
\end{itemize}
\end{KeyConcept}

\begin{ExampleBox}
\begin{itemize}
    \item \textbf{E-gold / paper gold} replacing physical gold (people still resist); \textbf{sex-determination technology} (people resist in the name of ethics); the \textbf{three-parent baby} (resisted in the name of ethics, values, religion); lal-tenn protest [as stated in lecture].
\end{itemize}
\end{ExampleBox}

\begin{QuickRevision}
\begin{itemize}
    \item Ogburn: material (technology) changes fast, non-material (values) slowly $\rightarrow$ cultural lag $\rightarrow$ strain; Luddites resist technology.
\end{itemize}
\end{QuickRevision}

\sectiondivider

\subsection{Conflict Theory (Marx)}

\begin{KeyConcept}
\begin{itemize}
    \item Society consists of \textbf{limited resources and opposing groups}; the opposition produces \textbf{conflict}, and conflict drives \textbf{social change}. The history of society is the history of \textbf{class struggle} (historical materialism): masters/slaves $\rightarrow$ feudalism (lords/serfs) $\rightarrow$ capitalism (bourgeoisie/proletariat) $\rightarrow$ socialism $\rightarrow$ communism. Conflict is inevitable; dialectics is the engine of history. (Marx: \textbf{polarisation} of classes; conflict between the two groups is inevitable.)
\end{itemize}
\end{KeyConcept}

\begin{QuickRevision}
\begin{itemize}
    \item Marx: limited resources + opposing classes $\rightarrow$ conflict $\rightarrow$ social change; history = class struggle (historical materialism).
\end{itemize}
\end{QuickRevision}

\sectiondivider

\section{Cyclical, Technological \& Modernization Theories}

\subsection{Answer-Writing: Hierarchy \& Exclusion as Impediments}

\begin{PYQLink}
\begin{itemize}
    \item \textbf{Keywords} to highlight: hierarchy, exclusion, transformation, impediment. Spend $\sim$20 seconds understanding the demand; \textbf{transformation = movement from closed to open society}.
    \item \textbf{Define hierarchy} = vertical arrangement of groups on \textbf{particular/ascriptive} standards (ascribed features) or \textbf{universal/meritocratic} standards (achievement orientation).
    \item \textbf{Define exclusion} = preventing members from full participation and realising their potential, on ascribed or achieved status.
    \item \textbf{Use two models}: (1) \textbf{Estate $\rightarrow$ Class} (France; the three estates were a closed system with hierarchy/exclusion justified by religion; the Enlightenment and French Revolution --- liberty, equality, fraternity, republic --- made it relatively open); (2) \textbf{Caste $\rightarrow$ Class} (India; Article 14 = equality of opportunity irrespective of caste/sex/race/residence $\rightarrow$ class society).
    \item \textbf{Neither is absolutely open nor closed} (IHDS: a specific group is over-represented in education/politics/occupation; caste has mobility --- Sanskritisation). In both open and closed systems hierarchy and exclusion persist --- but justified by \textbf{merit} (Davis \& Moore: inequality is functional) vs by \textbf{particular/ascriptive features}.
    \item \textbf{Tip}: do not answer on the basis of grand scholars (Marx, Weber) --- use \textbf{specialists of inequality} (Rousseau, André Béteille). A diagram can be drawn; end with the point that inequality is reproduced generationally (Béteille).
\end{itemize}
\end{PYQLink}

\begin{QuickRevision}
\begin{itemize}
    \item Define hierarchy/exclusion; two models (estate$\rightarrow$class, caste$\rightarrow$class); merit vs ascriptive justification; not absolutely open/closed; use specialists, not grand theorists.
\end{itemize}
\end{QuickRevision}

\sectiondivider

\subsection{Life Chances versus Lifestyle (PYQ)}

\begin{PYQLink}
\begin{itemize}
    \item The examiner is indirectly asking for the difference between \textbf{class and status group}. \textbf{Life chance} = the opportunities at an individual's disposal \textbf{owing to their market situation} (Weber's class definition: similar market situation $\rightarrow$ similar life chances). \textbf{Lifestyle} = one's \textbf{taste} --- the clothes you wear, the food you eat, the people you mingle with --- given by \textbf{status position}.
    \item Other PYQs the teacher flagged in this unit: ``\textbf{critically assess social mobility in closed and open systems}''; ``Davis \& Moore said social stratification is a \textbf{functional necessity and an unconscious device}''; ``\textbf{compare and contrast the contributions of Marx and Weber on stratification}''.
\end{itemize}
\end{PYQLink}

\begin{KeyConcept}
\begin{itemize}
    \item \textbf{Life chances come from class position} (wealth $\rightarrow$ better education, health, peer group, social capital); \textbf{lifestyle comes from status position}. Life chances $\rightarrow$ \textbf{Bourdieu's economic capital}; lifestyle $\rightarrow$ \textbf{symbolic capital}.
    \item \textbf{Example}: during lockdown, a person with a good market situation (smartphone, internet, urban location, education) could order groceries from Big Bazaar/Grofers and \textbf{practise social distancing with ease} --- while those without these life chances (migrants) had to step out and largely bore the COVID burden. A \textbf{nouveau-riche crorepati} who gambles and lacks taste has wealth but not lifestyle, whereas a \textbf{princely/billionaire family} carries a unique, distinctive lifestyle.
\end{itemize}
\end{KeyConcept}

\begin{QuickRevision}
\begin{itemize}
    \item Life chance (class/market situation $\rightarrow$ opportunities; economic capital) vs lifestyle (status $\rightarrow$ taste; symbolic capital); lockdown/grocer vs migrant example; nouveau riche vs princely family.
\end{itemize}
\end{QuickRevision}

\sectiondivider

\subsection{Conflict Theory: Ralph Dahrendorf}

\begin{KeyConcept}
\begin{itemize}
    \item \textbf{Dahrendorf} agrees with Marx that conflict exists, but says it is \textbf{not only between economic groups over economic issues} --- it can be between \textbf{non-economic groups over non-economic issues}, and such conflict also drives social change.
    \item Examples: the \textbf{two-nation theory} (Hindus vs Muslims --- Gandhi vs Jinnah); the \textbf{Moplah rebellion} (largely Hindu--Muslim, not economic); \textbf{Afghanistan} (Taliban = traditional vs modern --- a clash of civilisations); and \textbf{men vs women} (history = sex-class struggle; the \textbf{\#MeToo movement} is a gender conflict over dignity, not economics).
\end{itemize}
\end{KeyConcept}

\begin{QuickRevision}
\begin{itemize}
    \item Dahrendorf: conflict need not be economic (two-nation, Moplah, Afghanistan, \#MeToo).
\end{itemize}
\end{QuickRevision}

\sectiondivider

\subsection{Cyclical Theory (Spengler, Pareto, Sorokin)}

\begin{KeyConcept}
\begin{itemize}
    \item Every society, like an individual, is \textbf{born, matures, decays and dies} --- so it does not progress; after maturity it declines. Examples: the Roman Empire decayed; the Indus Valley civilisation decayed; the Congress system matured and then decayed (coalition politics = the beginning of its end).
\end{itemize}
\end{KeyConcept}

\begin{KeyConcept}
\begin{itemize}
    \item \textbf{Pareto's circulation of elites} (lions are replaced by foxes): elites are born, mature and decay; governing elite $\leftrightarrow$ non-governing elite circulate. In India this circulation happens \textbf{within} the governing elite (state $\leftrightarrow$ centre), not from outside --- only those with a political lineage survive (Kamal Haasan, from the South film industry, is an exception, not Bollywood).
\end{itemize}
\end{KeyConcept}

\begin{KeyConcept}
\begin{itemize}
    \item Every society oscillates between three belief systems: \textbf{Ideational} (religious knowledge), \textbf{Sensate} (scientific/sensory --- no belief in God), and \textbf{Idealistic} (a mix --- the point every society wants but never reaches). When a society moves back from the scientific to the religious, it shows \textbf{decline}.
    \item \textbf{Afghanistan} fits perfectly: it became modernised with American/NATO troops (modernising agents, education, girls' education, missionary missions) and then declined --- a living example of pendulum motion. So society does \textbf{not} inevitably progress.
\end{itemize}
\end{KeyConcept}

\begin{QuickRevision}
\begin{itemize}
    \item Spengler (born$\rightarrow$mature$\rightarrow$decay$\rightarrow$die); Pareto (circulation of elites, lions$\rightarrow$foxes); Sorokin (pendulum: ideational $\leftrightarrow$ sensate $\leftrightarrow$ idealistic; back to religion = decline; Afghanistan).
\end{itemize}
\end{QuickRevision}

\sectiondivider

\subsection{Technological Theory of Social Change (Veblen, Toffler)}

\begin{KeyConcept}
\begin{itemize}
    \item \textbf{Veblen} (inspired by Marx --- Marx first said technology drives social change) is a \textbf{technological determinist}: technology is the larger cause of social change. Technology $\rightarrow$ new means of production $\rightarrow$ relations of production $\rightarrow$ new mode of production.
    \item Steam engine $\rightarrow$ mass production $\rightarrow$ mass consumption $\rightarrow$ class structure (proletariat vs those who consume). Service sector $\rightarrow$ white-collar class. \textbf{Industry 4.0 / AI} $\rightarrow$ reserve army of labour, unemployment. The \textbf{Green Revolution} (high-yielding variety seeds + heavy machinery) led to the \textbf{forceful domestication of women} --- machinery legitimised man as breadwinner and eliminated women from agriculture.
    \item Technology has \textbf{positive and negative} sides.
\end{itemize}
\end{KeyConcept}

\begin{KeyConcept}
\begin{itemize}
    \item \textbf{Alvin Toffler} wrote \textit{Future Shock}, predicting two things: \textbf{information overload} (too much choice $\rightarrow$ confusion, ``madness'' --- Michel Foucault's \textit{Madness and Civilization}) and \textbf{death of permanence} (predicted in 1960: nothing will be permanent --- naukri.com means no permanent job, divorce, no ownership --- Uber instead of owning a car, renting houses, ``ghar ka khana'' dying).
    \item Negative side of technology: nuclear war, arms race (1950 onwards), trade war, global warming, ozone depletion, coronavirus (from a laboratory).
\end{itemize}
\end{KeyConcept}

\begin{QuickRevision}
\begin{itemize}
    \item Veblen (technological determinist): technology drives change (steam engine $\rightarrow$ mass production $\rightarrow$ class structure; Green Revolution $\rightarrow$ domestication of women).
    \item Toffler: information overload + death of permanence; negatives: arms race, global warming, COVID.
\end{itemize}
\end{QuickRevision}

\sectiondivider

\subsection{Modernization Theory: Founders \& Meaning}

\begin{KeyConcept}
\begin{itemize}
    \item Modernization has \textbf{no one meaning} and is \textbf{value-loaded} --- we equate ``modern'' with English, Western dress, Italian/Chinese food, a big car, and thereby call others ``non-modern,'' creating inequality.
    \item \textbf{Dipankar Gupta} (\textit{Mistaken Modernity}): the Indian middle class thinks it is modern because of consumption, but \textbf{modernity lies in values} --- it still practises caste endogamy and vote-bank politics.
    \item \textbf{Founders}: \textbf{Max Weber} (rationalisation --- replacing traditional/religious beliefs with the efficient, calculable, predictable); \textbf{Talcott Parsons} (pattern variables A $\rightarrow$ B, and \textbf{structural differentiation} --- specialised institutions: family once did everything; now school, state, market, law are specialised --- just like \textbf{Montesquieu's separation of powers}, each institution is a specialist in its own domain); \textbf{Karl Marx} (communism = a path to development; a classless society with no private property).
\end{itemize}
\end{KeyConcept}

\begin{KeyConcept}
\begin{itemize}
    \item Modernization theory developed in the \textbf{1950s} during the Cold War (1940--1990): should the world follow the capitalist bloc (America/Europe, Weberian/Parsonian) or the Soviet bloc (Marxist)? The theory explained \textbf{the poverty of Asian countries as the result of their own internal values} --- it called them ``traditional'' and offered Western help.
\end{itemize}
\end{KeyConcept}

\begin{QuickRevision}
\begin{itemize}
    \item Modernization = value-loaded, no single meaning; Gupta's Mistaken Modernity (values not consumption).
    \item Founders: Weber (rationalisation), Parsons (pattern variables + structural differentiation), Marx (communism).
    \item 1950s, Cold War; Asian poverty blamed on internal values.
\end{itemize}
\end{QuickRevision}

\sectiondivider

\subsection{Traditional versus Modern Society}

\begin{ComparisonBox}
\begin{tabularx}{\linewidth}{>{\bfseries}lYY}
\rowcolor{AccentDark}
\thead{Dimension} & \thead{Traditional society} & \thead{Modern society} \\
\midrule
Culture & Culture of \textbf{fatalism} (leave it to fate) & Culture of \textbf{innovation} \\
Family & \textbf{Large/joint families} $\rightarrow$ no savings & \textbf{Small families} (divorce, ``moving on'') \\
Economy & Lack of savings \& investment & Invest, not consume (Protestant work ethic; ``Hinduism = orthodox'') \\
State \& market & Government intervenes (MSP, minimum wage, petrol price) --- state and market are twins & Free market / neoliberal economy; state distant from market and religion \\
\bottomrule
\end{tabularx}
\end{ComparisonBox}

\begin{CriticalPoint}
\begin{itemize}
    \item Because the state regulates the market (MSP, minimum wage, petrol prices), India is \textbf{not a free-market nation-state} --- structural differentiation (a free market, a sovereign state) has not fully occurred.
\end{itemize}
\end{CriticalPoint}

\begin{QuickRevision}
\begin{itemize}
    \item Traditional: fatalism, large families, no saving, state intervenes; Modern: innovation, small families, invest, free market/secular state.
\end{itemize}
\end{QuickRevision}

\sectiondivider

\subsection{Rostow's Stages of Growth \& Daniel Lerner}

\begin{KeyConcept}
\begin{itemize}
    \item \textbf{Rostow} (financial advisor to the American President in the 1950s) wrote \textit{The Stages of Economic Growth: A Non-Communist Manifesto} --- the \textbf{aeroplane model} of modernization, with \textbf{five stages}: (1) traditional society, (2) preconditions for takeoff (jettison traditional/fatalist values, start saving, invest), (3) takeoff, (4) drive to maturity, (5) age of mass consumption.
    \item With Western money and advice, the ``aeroplane of economic growth'' would \textbf{taxi down the runway, pick up speed and become airborne} --- reaching a high standard of living and high income. This was propaganda to \textbf{Americanise the world}.
\end{itemize}
\end{KeyConcept}

\begin{KeyConcept}
\begin{itemize}
    \item \textbf{Daniel Lerner} studied the Middle East (Iran, Iraq, Jordan, Lebanon, Turkey) and called it underdeveloped for not following the Western model. For Lerner, \textbf{media is the driver of development} --- exposure through media (Afghanistan today; the Arab Spring) makes people reassess their lives and demand rights.
    \item If the \textbf{state captures media}, there is no structural differentiation --- media becomes the state's spokesperson (``godi'', not ``modi'') and society is not modern.
\end{itemize}
\end{KeyConcept}

\begin{QuickRevision}
\begin{itemize}
    \item Rostow: 5 stages (aeroplane model); Western money/advice $\rightarrow$ takeoff $\rightarrow$ mass consumption.
    \item Lerner: media = driver of development; captured media = not modern.
\end{itemize}
\end{QuickRevision}

\sectiondivider

\subsection{Critique of Modernization Theory}

\begin{CriticalPoint}
\begin{itemize}
    \item \textbf{1. It promotes Americanization/Europeanization} --- a propaganda for Western hegemony (instead of improving life, it standardises it; we become dependent on America).
    \item \textbf{2. It assumes unilinear development} --- in the name of evolution, \textbf{colonialism was justified}; today colonialism survives as \textbf{neo-colonialism} = \textbf{globalization} (McDonaldization, Coca-colonization).
    \item \textbf{3. It ignores ecological issues} --- the Western model of growth has damaged the global environment (hence the sustainable-goals debate).
    \item \textbf{4. It depicts the third world as ``third world''} --- an Orientalist picture of traditional society as a land of magic/superstition (from the Western lens vs the Nehru/Oriental-nationalist lens).
    \item \textbf{Hint to the next theory}: modernization theory says India is poor due to internal values; \textbf{dependency theory says India is poor because of Britain} --- the drain of wealth made India poor (``sone ki chidiya''). Value systems are only \textit{different}, not superior/inferior.
\end{itemize}
\end{CriticalPoint}

\begin{QuickRevision}
\begin{itemize}
    \item Critiques: Americanization; unilinear (justifies colonialism $\rightarrow$ neo-colonialism); ignores ecology; Orientalist ``third world'' framing.
    \item Dependency theory counters: external (drain of wealth), not internal values.
\end{itemize}
\end{QuickRevision}

\sectiondivider

\section{Dependency, World System, Development \& Globalization}

\subsection{Dependency Theory (A. G. Frank)}

\begin{KeyConcept}
\begin{itemize}
    \item \textbf{Dependency theory} (1950s--60s, founder \textbf{A. G. Frank}, a Marxist, case study \textbf{Latin America}) was a response to modernization theory. The cause of underdevelopment is \textbf{not cultural values but the dependency/exploitation} established between the West and the non-West.
    \item Frank divides the world into \textbf{metropolis} (developed --- the US, UK) and \textbf{satellites} (Latin America --- Peru, Venezuela, Bolivia). Satellites supply \textbf{raw materials} (petroleum/crude, copper, iron, spices, cotton, indigo); the metropolis \textbf{processes them into finished goods} and sells them back --- so wealth is drained (the trade always weighs in favour of the West). This is the \textbf{North--South divide} (the ``superior'' North).
\end{itemize}
\end{KeyConcept}

\begin{KeyConcept}
\begin{itemize}
    \item Colonialism divided the world into high- and low-income countries. After colonialism, the metropolis was \textbf{no longer a political coloniser but an economic coloniser} --- today this is called \textbf{globalization = neo-colonialism}. (China was a satellite in the 1950s, going Maoist, and is outside this model.)
    \item The \textbf{``drain of wealth''} idea actually predates Frank: \textbf{Dadabhai Naoroji} (who theorised that the drain of wealth made India --- the ``sone ki chidiya'' --- poor) would have been a \textbf{dependency theorist} if he were alive today. The \textbf{IMF's stated objective was to fix the balance-of-payment deficit} in traditional societies, which imported more than they could productively export; the \textbf{World Bank and IMF were a ``holy alliance''} set up to aid the non-Western societies.
    \item Note: a further flaw modernisation theory found in traditional societies was \textbf{corruption} --- when the state regulates the market, lobbies pressurise the government and corruption follows.
\end{itemize}
\end{KeyConcept}

\begin{KeyConcept}
\begin{itemize}
    \item Frank's solutions: \textbf{delinking} (overthrow the dependency relation, become self-sufficient through import substitution/in-house production --- from Marxist self-sufficiency) and \textbf{socialist revolution} (China, Cuba). The \textbf{Non-Aligned Movement} can be seen as an example of delinking (though different, as NAM is political).
    \item \textbf{Conclusion: development is a zero-sum game} --- it favours the West at the cost of the non-West. \textbf{``Development of underdevelopment''}: the development of the developed created the underdevelopment of the underdeveloped. \textbf{Ivan Illich} (who calls school an ``advertisement agency'') calls development a \textbf{``planned poverty''}.
\end{itemize}
\end{KeyConcept}

\begin{QuickRevision}
\begin{itemize}
    \item A. G. Frank: metropolis vs satellites; satellites supply raw materials, metropolis processes $\rightarrow$ drain of wealth; North-South divide.
    \item Post-colonial economic coloniser = neo-colonialism/globalization; solutions: delinking (import substitution), socialist revolution; development = zero-sum; Ivan Illich (planned poverty).
\end{itemize}
\end{QuickRevision}

\sectiondivider

\subsection{World System Theory (Immanuel Wallerstein)}

\begin{KeyConcept}
\begin{itemize}
    \item \textbf{Immanuel Wallerstein} (also Marxist, contemporary) divides the world into \textbf{core, semi-periphery and periphery} --- like a city's posh centre vs its edges.
    \item \textbf{Core} (US, Germany, France, UK): technologically/militarily/economically advanced; \textbf{exploits the periphery} through cheap imports, cheap labour, agricultural products and raw materials; \textbf{monopolises manufactured goods}.
    \item \textbf{Periphery} (Africa, Asia): provides \textbf{cheap raw materials including labour power} (Indian white-collar workers in Silicon Valley as cheap labour) --- its function is to ensure the core's productivity so the core expands.
    \item \textbf{Semi-periphery} (Russia, China, and today India): shares features of both (like Eric Olin Wright's contradictory class location) --- it benefits from unequal trade with the periphery but is disadvantaged in trade with the core; periphery is exploited from both sides.
\end{itemize}
\end{KeyConcept}

\begin{CriticalPoint}
\begin{itemize}
    \item When the periphery tries to rise, the core \textbf{resists} --- the US imposed sanctions on China (tariff and \textbf{non-tariff/phytosanitary barriers} = weapons to maintain exclusion). India has developed bargaining power (20\% local-outsourcing clauses; refusing \textbf{RCEP}; Make in India, local-for-local).
\end{itemize}
\end{CriticalPoint}

\begin{ComparisonBox}
\begin{itemize}
    \item Both are \textbf{Marxist} and both see development of one at the cost of another. \textbf{Use A. G. Frank (metropolis--satellite) for questions on colonialism; use Wallerstein (core--periphery) for questions on contemporary globalization.}
\end{itemize}
\end{ComparisonBox}

\begin{QuickRevision}
\begin{itemize}
    \item Wallerstein: core (US/UK/Germany/France) exploits periphery (Africa/Asia) + semi-periphery (Russia/China/India).
    \item Periphery supplies raw material incl. labour power; core resists rise via tariff/non-tariff barriers; Frank for colonialism, Wallerstein for globalization.
\end{itemize}
\end{QuickRevision}

\sectiondivider

\subsection{Critique of Dependency Theory}

\begin{CriticalPoint}
\begin{itemize}
    \item \textbf{1. Colonialism also benefited} --- railways, modern education; the \textbf{poorest countries are those never colonised} (e.g., \textbf{Ethiopia}).
    \item \textbf{2. It blames only external causes and ignores internal ones} --- ethnic/communal conflict, caste conflict, gender inequality, corruption, and \textbf{landlocked geography} (Chhattisgarh, Jharkhand, Bihar, eastern Maharashtra/Orissa are landlocked and underdeveloped).
    \item \textbf{3. Its solutions (import substitution, self-sufficiency, socialist revolution) have miserably failed} --- socialist revolution produced totalitarianism/genocide, and poverty reproduced.
    \item \textbf{4. It portrays all satellites as one} --- it treats all low-income/least-developed countries as the same.
\end{itemize}
\end{CriticalPoint}

\begin{QuickRevision}
\begin{itemize}
    \item Critiques: colonialism benefited (Ethiopia); ignores internal causes (conflict, landlocked); solutions failed; treats all satellites as one.
\end{itemize}
\end{QuickRevision}

\sectiondivider

\subsection{Development: Meaning \& Evolution}

\begin{KeyConcept}
\begin{itemize}
    \item Till the 1990s, \textbf{GDP/GNP was the measure of development} --- development was \textbf{material} (a materialist approach). Post-1990s came \textbf{HDI}; post-2000 \textbf{GNH (Gross National Happiness)} --- happiness as development (\textbf{Bhutanese/buddhist economy}; Bhutan has 70\% protected forest and rejects McDonald's/Nike, denying even India's transport corridor on ecological grounds).
    \item So \textbf{growth $\neq$ development}: growth is economic, development is human/non-material. \textbf{Amartya Sen's capability approach}; \textbf{green GDP}; \textbf{MDG (2000--2015)} and \textbf{SDG (2015--2030)}.
    \item The \textbf{discourse of development has evolved} --- from ``development of growth'' (material) to ``growth of development'' (human). Human development is a \textbf{utopia} --- a reference point we work towards.
\end{itemize}
\end{KeyConcept}

\begin{ComparisonBox}
\begin{tabularx}{\linewidth}{>{\bfseries}lYY}
\rowcolor{AccentDark}
\thead{Dimension} & \thead{Mainstream (growth)} & \thead{Alternative (social transformation)} \\
\midrule
Objective & Accumulation of wealth & Capacitation \& human development \\
Resources & Capital, technology, trade, foreign investment & Human skills, local resources, social \& human capital, local knowledge \\
Agency & Market-led \& state-led & Community, NGOs, civil society, local government \\
Method & Industrialisation, export-led growth, innovation & Participation, micro-credit, sustainability, democratisation \\
Indicator & GDP, GNP & Green GDP, SDG, HDI, GNH \\
\bottomrule
\end{tabularx}
\end{ComparisonBox}

\begin{ExampleBox}
\begin{itemize}
    \item \textbf{People's participation}: tribal-led forest conservation (local knowledge), the \textbf{People's Biodiversity Register (PBR)}, \textbf{Joint Forest Management}, Panchayati Raj, the Ministry of Cooperation/cooperatives, SHGs, micro-credit. Development was \textbf{top-down, now bottom-up} --- this is \textbf{ethno-development} (Hetne).
\end{itemize}
\end{ExampleBox}

\begin{QuickRevision}
\begin{itemize}
    \item GDP/GNP (material, till 90s) $\rightarrow$ HDI $\rightarrow$ GNH/Bhutan (happiness); growth $\neq$ development; Sen's capability approach; MDG$\rightarrow$SDG; green GDP.
    \item Mainstream (growth/accumulation/market) vs alternative (transformation/capacitation/community); ethno-development = bottom-up.
\end{itemize}
\end{QuickRevision}

\sectiondivider

\subsection{Globalization: Dimensions \& Schools}

\begin{KeyConcept}
\begin{itemize}
    \item Globalization = the \textbf{flow of objects, ideas and people freely, without regulation}. It is not new (ancient --- Hiuen Tsang studying Buddhist literature, Marco Polo, the Greek invasion), but today it is \textbf{rapid and inclusive}.
    \item The \textbf{global citizen}: an American drives a car \textbf{made in Germany} (steel from Korea, tyres from Malaysia, gasoline from UK/BP pumped from Africa and shipped by a Greek company assembled in France), sips \textbf{Starbucks} coffee (beans from Africa/Brazil/Peru, processed in China), and watches the \textbf{Afghanistan crisis} on a mini-TV --- all at once.
    \item Globalization of \textbf{ideas} is as real as of people: liberty, equality, fraternity came from France; \textbf{nationalism} from France; \textbf{totalitarianism} from Europe; \textbf{feminism and the \#MeToo movement} from Europe/America into India.
\end{itemize}
\end{KeyConcept}

\begin{KeyConcept}
\begin{itemize}
    \item \textbf{Liberalization} (minimum regulation, free market, deregulation); \textbf{modernization/westernization} (replacement of traditional by modern ideas --- nationalism from France, totalitarianism from Europe); \textbf{internationalization} (cross-border relations --- diplomacy, the World Economic Forum); \textbf{universalization} (spread of products/ideas/services); \textbf{de-territorialization} (moving beyond territories --- the fifth sums up the other four).
\end{itemize}
\end{KeyConcept}

\begin{KeyConcept}
\begin{itemize}
    \item \textbf{Hyper-globalizers}: \textit{pro-globalist} (optimist) --- \textbf{Thomas Friedman} (``neoliberalism is a golden straitjacket''; trickle-down, like Adam Smith's invisible hand; the World Bank \& IMF as the ``holy alliance''); \textbf{Kenichi Ohmae} (``borderless world''); \textbf{Marshall McLuhan} (``global village'' --- instant information, informed cosmopolitan citizens). \textit{Anti-globalist} (pessimist) --- globalization is \textbf{cultural imperialism}, homogenisation, Americanization. Ohmae's ``borderless world'' shades into \textbf{denationalization} --- the nation-state losing relevance (multilateral bodies judge our ``internal'' disputes; the \textbf{Foreign Investment Promotion Board was abolished} and FDI moved to the \textbf{automatic route} toward 100\%, so the state becomes a mere facilitator for global corporations).
    \item \textbf{Skeptics} (doubtful, not anti): \textbf{Paul Hirst} --- what we call globalization is really \textbf{regionalization} (NAFTA, EU, African Union, SAARC, ASEAN; the G7/G8 are closed, exclusive groupings). Even the UN Security Council is an \textbf{exclusionary social closure}: India contributes the \textbf{largest UN peacekeeping force} yet is denied a permanent seat, because the existing members do not want outsiders in.
    \item \textbf{Transformationalists}: globalization is a \textbf{two-way process} that runs on the \textbf{adaptive capacity of nation-states} (vaccine nationalism; India reverting to government-approval route for Chinese investment) --- the nation-state never loses relevance, it adapts. \textbf{Roland Robertson} = \textbf{glocalization}; \textbf{Anthony Giddens} = people become more \textbf{``glocal''} (hybrid identities --- McAloo Tikki, Domino's paneer tikka, kurta with jeans, Panchayati Raj = democracy + village assembly). Globalization is therefore \textbf{not a threat to indigenous culture} --- it produces \textbf{hybridization}.
\end{itemize}
\end{KeyConcept}

\begin{QuickRevision}
\begin{itemize}
    \item Globalization = free flow of objects/ideas/people; five dimensions (Scholte): liberalization, westernization, internationalization, universalization, de-territorialization.
    \item David Held's schools: hyper-globalizers (Friedman/Ohmae/McLuhan vs cultural imperialism), skeptics (Hirst --- regionalization), transformationalists (two-way, adaptive states; glocalization --- Robertson/Giddens).
\end{itemize}
\end{QuickRevision}

\sectiondivider

\section{Work \& Economy}

\subsection{Why Sociology Studies Work (Durkheim, Marx, Weber)}

\begin{KeyConcept}
\begin{itemize}
    \item Work seems like an economist's domain, but \textbf{work is socially constructed} --- not everything we do is ``work'' (work = the paid job; a housewife's work is unpaid and not counted).
    \item \textbf{Durkheim} (\textit{The Division of Labour}, 1890s): division of labour is a \textbf{moral force}, but it took an \textbf{abnormal form} (anomic, pathological, inadequately coordinated) during the \textbf{anomic phase} after the Industrial Revolution; with the state, registered/cooperative bodies and associations, the normal division of labour will be restored.
    \item \textbf{Marx}: work in capitalism is a \textbf{site of alienation} --- the kind of work, the wage, and the hours are not chosen by the worker.
    \item \textbf{Weber}: work takes a \textbf{bureaucratised form} --- recruitment by merit and written rules, coordination by hierarchy (chain of command) and a code of conduct.
\end{itemize}
\end{KeyConcept}

\begin{QuickRevision}
\begin{itemize}
    \item Work is socially constructed (paid vs unpaid). Durkheim (abnormal division of labour in the anomic phase); Marx (alienation); Weber (bureaucratisation).
\end{itemize}
\end{QuickRevision}

\sectiondivider

\subsection{Paid \& Unpaid Work}

\begin{KeyConcept}
\begin{itemize}
    \item \textbf{Paid and unpaid are both work}, but only paid work is \textbf{registered in GDP}. Within paid work there is the formal/informal distinction. Defining work as ``paid'' automatically \textbf{excludes a large chunk} of the population --- the idea of work is exclusionary.
\end{itemize}
\end{KeyConcept}

\begin{QuickRevision}
\begin{itemize}
    \item Paid + unpaid both work; only paid in GDP; ``work = paid'' is exclusionary.
\end{itemize}
\end{QuickRevision}

\sectiondivider

\subsection{Formal \& Informal Sector (Keith Hart)}

\begin{KeyConcept}
\begin{itemize}
    \item \textbf{Formal/informal are only models (ideal types)} devised by sociologists/economists --- they do not exist sharply in reality (a parking attendant with a slip is ambiguous).
    \item \textbf{Informal sector} features: \textbf{self-employed} (freelancer or own business), \textbf{wage-employed}, \textbf{low-skilled}, \textbf{ease of entry}, \textbf{micro-to-small scale}, \textbf{no restriction on hours/minimum wage} (no social security).
    \item \textbf{Formal sector} features: \textbf{high-skilled}, \textbf{entry by merit/rules}, \textbf{medium-to-large scale}, \textbf{predefined hours/wage + social security} (provident fund, insurance).
    \item \textbf{Informal employment vs informal economy} (the labour sociologist \textbf{Sharad Bhaumik} is associated with this point): the informal economy is \textbf{unregistered} (shadow/grey economy, not paying full tax); \textbf{informal employment exists even inside the formal sector} (the DU Sociology Xerox shop, the CP parking association's vendors, an independent advocate, per-hour/share teachers in Rajinder Nagar).
    \item Formal/informal is decided by \textbf{whether transactions are recorded with the government} --- a government officer's \textbf{stock-market investments} (SEBI-registered, demat account, security transaction tax) are formal, but the same person \textbf{betting on a lottery/gambling app} is informal.
\end{itemize}
\end{KeyConcept}

\begin{ExampleBox}
\begin{itemize}
    \item \textbf{Girdhari Lal} in \textbf{Ambience Mall} = formal (pays tax), but his two ``chotus'' (helpers) are informal; the \textbf{ORN paan shop} = registered but its helpers informal; the \textbf{under-the-tree paan shop} = unregistered, wholly informal. Every time you ride an auto-rickshaw, you are in the informal economy. \textbf{Formal--informal is a continuum} in which members move to and fro.
\end{itemize}
\end{ExampleBox}

\begin{QuickRevision}
\begin{itemize}
    \item Keith Hart: formal/informal = ideal types. Informal: self-employed, low-skilled, easy entry, small scale, no security. Formal: merit entry, large scale, social security.
    \item Informal employment inside formal sector; formal--informal = continuum.
\end{itemize}
\end{QuickRevision}

\sectiondivider

\subsection{Informalization of Work}

\begin{KeyConcept}
\begin{itemize}
    \item Per the \textbf{ILO}: till the 1980s--90s there was a large formal sector and small informal hubs assisting it (McDonald's is formal, but its contractual/part-time workers and outsourced potato farmers are informal --- a \textbf{global value chain}). From 2000 onwards, \textbf{formal is contracting and informal expanding} --- the \textbf{informalization of work} (accelerated by the rise of the \textbf{service sector}).
\end{itemize}
\end{KeyConcept}

\begin{KeyConcept}
\begin{itemize}
    \item \textbf{1. Flexibilization of work} --- contractual workers, dependent employees, self-employed; the Uber driver is a \textbf{camouflaged wage worker} (Jan Breman) --- ``be your own boss'' but the car is leased; women seek part-time/flexible jobs.
    \item \textbf{2. Globalization} --- global value chains; cheap labour from Asia; shift from \textbf{secure to precarious self-employment}.
    \item \textbf{3. Capitalism} --- the \textbf{law of capital accumulation} (profit $\rightarrow$ cut expenditure $\rightarrow$ fire workers $\rightarrow$ \textbf{industrial reserve army} = a pool of informal labour); \textbf{flexibilization of production} (contracting out; \textbf{piece-rate} instead of time-wage; casual instead of permanent).
    \item \textbf{4. Introduction of technology} --- replaces low-skilled jobs (typist, calculator); Uber threatened rickshaw/auto drivers; women (in low-skilled jobs) are the biggest victims; \textbf{capital-intensive growth} puts workers in a race with technology and fellow workers to remain relevant.
    \item \textbf{5. Demographic factors} --- surplus labour, rural-urban migration, urbanisation: migrants land first in the \textbf{informal economy}.
    \item \textbf{6. Subcontracting of work} --- contracting work out in a long chain brings informality into the formal sector.
\end{itemize}
\end{KeyConcept}

\begin{QuickRevision}
\begin{itemize}
    \item ILO pattern: formal contracting, informal expanding $\rightarrow$ informalization (service sector).
    \item Factors: flexibilization (camouflaged wage workers), globalization (cheap labour, precarious self-employment), capitalism (capital accumulation $\rightarrow$ reserve army; piece-rate), technology (replaces low-skilled, women), demographic (migration $\rightarrow$ informal), subcontracting.
\end{itemize}
\end{QuickRevision}

\sectiondivider

\subsection{Precariat Class \& Foot-loose Labour}

\begin{KeyConcept}
\begin{itemize}
    \item \textbf{Guy Standing} (book, 2013) coined the \textbf{precariat class} --- a new class from flexibilization + globalization, not permanently positioned in any organisation, fully freelancing (Uber/Swiggy/Amazon delivery boys --- the \textbf{platform economy}). ``Precarious'' = dangerous/vulnerable; their tomorrow is never certain.
\end{itemize}
\end{KeyConcept}

\begin{KeyConcept}
\begin{itemize}
    \item Migrant labourers act as \textbf{foot-loose labour} (Jan Breman) --- freelancers moving to and fro, vulnerable to forced labour. In lockdown they were the biggest victims of labour-law suspension --- \textbf{circular migration} (rural $\rightarrow$ urban $\rightarrow$ back). They die on the way home because they have no savings (the informal sector has no saving, only consumption).
\end{itemize}
\end{KeyConcept}

\begin{QuickRevision}
\begin{itemize}
    \item Precariat class (Guy Standing): platform economy (Uber/Swiggy/Amazon), precarious, insecure.
    \item Foot-loose labour (Jan Breman): circular migration, no savings, lockdown victims.
\end{itemize}
\end{QuickRevision}

\sectiondivider

\subsection{Adam Smith, Taylor \& Ford (Fordism)}

\begin{KeyConcept}
\begin{itemize}
    \item \textbf{Adam Smith}: \textbf{division of labour} = the source of productivity (break work into tasks; 10 men do in 2 minutes what one did in an hour) --- but this hides \textbf{informalization} (the minute tasks are done by low-skilled, replaceable freelancers).
    \item \textbf{Frederick Taylor}: \textbf{scientific management} --- supervisors/CCTV monitor workers; piece-work is timed.
    \item \textbf{Henry Ford}: production is not enough, \textbf{consumption is equally important} --- he \textbf{raised wages from \$2 to \$5 a day} (resolving the \textbf{inherent contradiction of capitalism that Marx exposed}: produce, but with nobody earning enough to consume), introduced the \textbf{assembly line} (an idea he saw in \textbf{slaughter houses/sweat shops} --- where others saw animal-rights violation he saw \textbf{productivity}, a ``phenomenology'' of perspective), and made the \textbf{T-model} (all black). \textbf{Fordism = mass production + mass consumption.}
\end{itemize}
\end{KeyConcept}

\begin{KeyConcept}
\begin{itemize}
    \item Mass production; \textbf{high unionisation}; \textbf{long-term commitment to workers}; \textbf{collective bargaining agreements} (firm--union agreement on conditions, wages, seniority, benefits). The trade unions themselves were \textbf{bureaucratised} (a leader, a chain of command, a subscription fee --- and \textbf{not all workers} were members).
    \item \textbf{Decline of Fordism}: it was \textbf{domestic-oriented} (could not serve global demand), produced \textbf{standardised goods} (no innovation --- just as \textbf{LG and Sony failed in smartphones} for lack of innovation while \textbf{Vivo/Oppo/Mi} out-innovated), could not cater to individual customers, and its goods were \textbf{expensive} (because high wages raised production cost) --- German/Japanese cars (cheaper, more efficient, colourful) beat the T-model.
\end{itemize}
\end{KeyConcept}

\begin{QuickRevision}
\begin{itemize}
    \item Adam Smith (division of labour $\rightarrow$ productivity + hidden informalization); Taylor (scientific management/CCTV); Ford (production + consumption, \$2$\rightarrow$\$5, assembly line from slaughter houses, T-model).
    \item Fordism = mass production + consumption; declined (domestic, standardised, expensive).
\end{itemize}
\end{QuickRevision}

\sectiondivider

\subsection{Post-Fordism}

\begin{ComparisonBox}
\begin{tabularx}{\linewidth}{>{\bfseries}lYY}
\rowcolor{AccentDark}
\thead{Dimension} & \thead{Fordism} & \thead{Post-Fordism} \\
\midrule
Production & Uniform at one location (rigid assembly line) & Flexible at distinct locations (computer-aided design, mass customisation) \\
Economy & Economies of scale (mass volume) & Economies of scope (variety/range) \\
Workforce & Rise of manufacturing workers & Rise of service \& white-collar workers \\
Goods & Mass production of homogeneous goods & Small batches of heterogeneous goods \\
Labour & Bureaucratised trade unions, collective bargaining & Decline of trade unions (individualism; workers compete) \\
\bottomrule
\end{tabularx}
\end{ComparisonBox}

\begin{ExampleBox}
\begin{itemize}
    \item \textbf{Dell}: earlier it made computers first then collected money; today it collects money first, then builds the customised (assembled) computer --- the era of customisation. A personalised ``I Love Nishat'' t-shirt is small-batch, not mass-produced.
\end{itemize}
\end{ExampleBox}

\begin{QuickRevision}
\begin{itemize}
    \item Post-Fordism: flexible production, economies of scope, service/white-collar workers, small heterogeneous batches, declining trade unions (individualism).
\end{itemize}
\end{QuickRevision}

\sectiondivider

\subsection{Feminization of Workforce}

\begin{KeyConcept}
\begin{itemize}
    \item \textbf{Feminization of workforce} = (1) replacement of the male workforce by female, or (2) the increased presence of women (not necessarily at men's cost --- feminists demand equal wage, no gender gap, equal respect, not that men leave). With rising feminism, sociologists observe the \textbf{``fading of blue and the coming of pink''} --- a formerly male job becomes a female one.
    \item \textbf{Feminization of work} = the gendering of work itself --- ``women's work'' (pink-collar jobs); work has a gender.
    \item \textbf{Reasons}: occupational segregation (horizontal + vertical), part-time work (women run the ``second shift''), and the \textbf{wage gap}.
\end{itemize}
\end{KeyConcept}

\begin{KeyConcept}
\begin{itemize}
    \item \textbf{Economists}: \textbf{compensating differentials} (women prefer comfortable conditions so are paid less --- treating her as a free agent) and the \textbf{crowding effect} (many women crowd low-number, low-prestige positions, so wages fall).
    \item \textbf{Sociologists}: \textbf{devaluation} (jobs are low-paid because they are seen as low-status ``women's jobs'') and \textbf{queuing} (better-paying jobs are hoarded by men, so women queue for poorly rewarded occupations).
\end{itemize}
\end{KeyConcept}

\begin{CriticalPoint}
\begin{itemize}
    \item \textbf{Heidi Hartmann} shows the link: \textbf{capitalism and patriarchy are twin brothers} --- low-paid, low-status women's work creates women's \textbf{dependency on men} (which is why marital alliances are patrilocal).
\end{itemize}
\end{CriticalPoint}

\begin{QuickRevision}
\begin{itemize}
    \item Feminization of workforce (replacement / increased presence) vs feminization of work (pink-collar); reasons: segregation, part-time/second shift, wage gap.
    \item Wage gap: economists (compensating differentials, crowding) vs sociologists (devaluation, queuing); Hartmann (capitalism + patriarchy = twins).
\end{itemize}
\end{QuickRevision}

\sectiondivider

\subsection{Female Labour Force Participation \& the Feminization-U}

\begin{DataFact}
\begin{itemize}
    \item \textbf{Female Labour Force Participation Rate (FLFPR)} in India $\approx$ \textbf{26\%} (dropped to $\sim$16\% during COVID per a student's correction); globally $\approx$ \textbf{45\%}. \textbf{Amartya Sen's ``missing women''}.
    \item The \textbf{Feminization-U hypothesis}: as GDP per capita rises, FLFPR first \textbf{declines} then \textbf{rises} (a U-curve).
\end{itemize}
\end{DataFact}

\begin{KeyConcept}
\begin{itemize}
    \item \textbf{Incompatibility of work and family} (women stop working after marriage; a working woman is a taboo in traditional families); \textbf{scarcity of suitable non-farm jobs near home}; \textbf{NREGA} (33\% women's reservation) made rural women \textbf{aspirational} --- they now look down on farm work; \textbf{stigma against women working outside}; \textbf{sexual violence / unsafe workplaces} (\textbf{public patriarchy} --- Muslim and Dalit women are the main victims); \textbf{lack of formal jobs}; \textbf{maternity} (maternity leave made the state patriarchal; now paternal leave is state-wise and short); and \textbf{education} (women voluntarily opt out to study --- Beti Bachao Beti Padhao).
    \item The drop is mainly \textbf{rural} --- rural women (engaged in agriculture) are leaving farm work for education/non-farm aspirations.
\end{itemize}
\end{KeyConcept}

\begin{CriticalPoint}
\begin{itemize}
    \item Technology is \textbf{gendered} (\textbf{femtech}): some technologies are ``women-only'' (the vacuum cleaner), some ``men-only'' (the smartphone is seen as a male technology; heavy machinery in agriculture is men-friendly). Heavy machinery replaced low-skilled women's jobs (winnowing, harvesting) in the Green Revolution, eliminating them from agriculture.
\end{itemize}
\end{CriticalPoint}

\begin{QuickRevision}
\begin{itemize}
    \item FLFPR India $\sim$26\% (global 45\%); Feminization-U (declines then rises with GDP).
    \item Reasons: work/family conflict, no non-farm jobs, NREGA aspirations, stigma, public patriarchy, maternity, education (voluntary); rural-led drop; gendered technology (femtech).
\end{itemize}
\end{QuickRevision}

\sectiondivider

\subsection{Feminization of Agriculture (Bina Agarwal)}

\begin{KeyConcept}
\begin{itemize}
    \item \textbf{Feminization of agriculture} may look like rising women's participation, but it is actually the result of \textbf{men's migration to urban areas} (agrarian distress = push factor; de-peasantization) --- women fill the \textbf{vacuum} left behind on the ``brown pastures.''
    \item \textbf{90 million} women are in the rural workforce; \textbf{60 million} are \textbf{agrarian labourers} (working on someone else's land, with no land of their own).
\end{itemize}
\end{KeyConcept}

\begin{DataFact}
\begin{itemize}
    \item \textbf{Bina Agarwal} notes: \textbf{56.4\% of rural households own no agricultural land}; women do \textbf{80\% of farm work but own only 13\% of land} (Oxfam data). Farmer associations and markets are controlled by men; SHG empowerment works \textbf{only within the group} --- at home the elderly man still decides.
    \item Therefore she prefers the term \textbf{de-feminization of agriculture} --- women are workers, not owners. This is \textbf{feminization of agrarian distress} $\rightarrow$ \textbf{feminization of poverty}.
\end{itemize}
\end{DataFact}

\begin{QuickRevision}
\begin{itemize}
    \item Feminization of agriculture = male migration + agrarian distress (push factor), not choice; 60m/90m women are agrarian labourers.
    \item Bina Agarwal: 56.4\% landless households; women 80\% work but 13\% land; SHG empowerment limited to the group; de-feminization of agriculture; feminization of poverty.
\end{itemize}
\end{QuickRevision}